\documentclass[12pt]{article}
\usepackage{graphicx}
\usepackage{amsmath}
\usepackage{url}
\usepackage{booktabs}
\usepackage[margin=1in]{geometry}
\usepackage{natbib}
\usepackage{xurl}
\usepackage{setspace}
\doublespacing

\title{Uneven Legalization, Spatial Demand Spillovers, and the Regulation of Entry:\\
Evidence from U.S.\ Cannabis Markets}
\author{Mahyar Rezazadeh \\ Babak Sahragard}
\date{August 2026}

\begin{document}
\maketitle

\begin{singlespace}
\begin{abstract}
\noindent Uneven state-level legalization of recreational cannabis creates a spatial demand spillover: counties
that share a border with a still-prohibited neighboring state face captive cross-border demand that interior
counties of the same state do not. We ask two questions about this spillover. First, does it produce a border
premium in legal cannabis markets. Second, does entry regulation determine which margin absorbs that premium
-- existing sellers' revenue, or the number of sellers and the market's total size. Using a county-year panel
spanning Colorado, Washington, Oregon, Montana, and New Mexico (2015--2026), with a border classification
defined by literal county-to-state-line geometric adjacency rather than a distance threshold, and time-varying
rather than a fixed snapshot -- accounting for New Mexico's 2022 legalization, Oregon's 2015 legalization, and
Oregon's own 2017--2018 transition relative to California and Nevada -- we find a robust border premium on
revenue per dispensary (+48\%, county-clustered $p=0.18$), revenue per capita (+173\%,
$p=0.001$), and dispensary density (+98\%, $p<0.001$). Whether entry is capped, using an institutional
classification that treats binding local-government control over licensing as equivalent to a state cap, does
not moderate the revenue-per-store margin, but it sharply moderates the other two: capped entry reduces the
border premium's multiplier on revenue per capita by 47\% in relative terms (from roughly 173\% to roughly
44\% in level terms, $p=0.048$) and on dispensary density by 43\% in relative terms ($p=0.009$). Oregon's own
within-state entry freeze in January 2022 produces the same-signed, if
underpowered, pattern. A dispensary-level panel built for New Mexico -- the one sample state where individual
firm revenue is observable historically -- confirms the border premium is not a county-aggregation artifact
(+141\% at the level of the individual dispensary, $p=0.001$, significant in every sample year). The results
support a specific mechanism -- entry regulation determines whether spatial
demand spillovers are absorbed by incumbent revenue or by market expansion -- while also reporting, without
forcing a cleaner story, that the mechanism does not extend to the intensive margin the way it was originally
hypothesized to. We directly tested whether Washington's 2016 cap increase should be treated as a shift toward
free entry and found it should not: Washington's general license application window closed permanently in
March 2016 and never reopened, and its post-2016 border counties do not behave like a genuinely free-entry
benchmark (New Mexico) on dispensary density, confirming the paper's classification of
Washington as capped throughout.

\vspace{0.5em}
\noindent\textbf{JEL codes:} L51, H71, R12, L81
\vspace{0.3em}

\noindent\textbf{Keywords:} cannabis, cross-border shopping, market entry, license caps, entry regulation
\end{abstract}
\end{singlespace}

\section{Introduction}
\label{sec:intro}

By the end of 2025, 24 U.S.\ states had legalized recreational cannabis, though Virginia had not yet
established an operating retail market, leaving 23 states with a functioning legal market and 27 states,
plus the remaining territory, where recreational retail sale remained prohibited. This patchwork did not
arrive all at once, and it has not stopped changing: a state that borders only prohibition-states in one year
can border a freshly legal neighbor the next. That timing creates a natural, if underused, source of
identifying variation. A county sitting near a still-prohibited state faces a captive pool of residents who
cannot buy cannabis legally at home; a county in the interior of the same legal state does not. If cannabis
demand is at all price- or access-sensitive across a state line -- which the existing literature on cannabis,
alcohol, and other regulated goods says it is -- that captive demand should show up somewhere in the legal
market's own data.

Where it shows up, however, is not obvious, and that is the question this paper is really about. A dispensary
facing extra demand from across the border can respond in at least two ways: it can sell more at roughly the
same price and margin, raising its own revenue, or new dispensaries can open nearby to meet the extra demand,
raising the county's total revenue and store count without necessarily raising revenue at any one store. Which
of these happens is not a fact about demand; it is a fact about supply-side regulation. A state that caps the
number of licenses, whether by an explicit numerical ceiling or by letting local governments block entry
outright, closes off the second channel by construction. A state that does not cap entry leaves it open. This
paper's central hypothesis is that cross-border demand creates a revenue premium in border counties, and that
entry regulation determines whether that premium is absorbed by existing sellers (revenue per store) or by
market expansion (revenue per capita and dispensary count) -- a specific, testable version of a standard
regulatory-economics prediction: a binding entry restriction forces a demand shock into incumbent rents rather
than into quantity.

We test this using a county-year panel covering five states -- Colorado, Washington, Oregon, Montana, and New
Mexico -- assembled from each state's own revenue-reporting agency (2014--2026, depending on the state's own
retail-launch date), together with an institutional classification of each state's entry-licensing regime built
from the actual statutes and regulations, not from a state's own label for itself. Two features of the design
depart from prior work in this space, including our own earlier work on Colorado and Washington alone. First,
a county's exposure to a nonlegal-state border is treated as time-varying: because several of the sample
states' neighbors themselves legalized cannabis partway through the sample window (New Mexico in 2022, Oregon
in 2015, and, more narrowly, California and Nevada relative to Oregon's own 2017 observations), a border
classification fixed at a single snapshot date misclassifies a meaningful share of county-years. We recompute
it year by year. Second, entry regulation is classified by where the binding constraint on licensing actually
sits, not by whether a state imposes an explicit numerical cap. Colorado has no state-level cap, but roughly
half its counties ban cannabis retail outright through local option, a restriction documented independently by
\citet{PayanBrownSong2021}; Montana has no numerical cap either, but its 2021 enabling statute requires
counties that voted against the state's 2020 legalization initiative to affirmatively opt in before a
dispensary can locate there. Both are coded as capped-entry states for this reason, a coding choice we explain
and defend in Section~\ref{sec:institutional} rather than assert.

We find, first, that cross-border demand generates a substantial and statistically robust premium in legal
cannabis markets, though not uniformly across every outcome we test: border counties earn 173 percent more
revenue per capita than interior counties of the same state ($p=0.001$) and have 98 percent more
dispensaries per capita ($p<0.001$), pooling across the five states. We do not find a robust border premium
on revenue per individual dispensary store ($p=0.181$), and report this directly rather than force a uniform
story across all three margins. The per-capita magnitude is consistent with \citet{HansenMillerWeber2020}'s
finding, in the mirror-image setting, that Washington retailers near the Oregon border lost roughly 36 percent
of their sales when Oregon's own retail market opened in October 2015 -- both results describe cross-border
demand effects of several tens of percentage points, not a marginal phenomenon.

Second, entry regulation determines whether this premium survives. Capping entry -- whether through a direct
state-level license ceiling (Washington), county-level opt-out bans (Colorado), or a county opt-in requirement
(Montana) -- reduces the border premium's multiplier on revenue per capita by 47 percent in relative terms
($p=0.048$) and on dispensary density by 43 percent in relative terms ($p=0.009$), similar magnitudes
that point to a single mechanism: capped entry blocks the quantity margin specifically, not the price margin.
Consistent with this, the interaction on revenue per store is a precise, economically negligible null ($+8.5$
percent, $p=0.799$) -- capped entry does not redirect the premium into higher per-store prices, it simply
prevents new stores from opening to capture it, consistent with \citet{HollenbeckUetake2021}'s finding that
Washington's own capped licensing regime supports meaningful incumbent markups rather than eliminating rents
at the store level.

Third, this has a real fiscal cost. Applying the estimated gap between the free-entry and capped-entry
revenue-per-capita premium to 2025 border-county revenue in the three states where the cap is unambiguously
binding (Colorado, Washington, Montana) implies these counties would need roughly \$470 million in additional
aggregate revenue to match what free entry would generate for counties of their size, and, at each state's own
statutory cannabis tax rate, roughly \$107 million in foregone state tax revenue in a single year
(Section~\ref{sec:magnitude}). This is considerably larger than \citet{Thomas2019LicenseQuotas}'s structural
estimate that free entry would have raised total surplus by roughly 18 percent under Washington's own license
quota -- an expected difference, since a revenue-gap accounting includes normal costs and rents that a surplus
calculation nets out, but the two estimates agree qualitatively that Washington's entry cap imposes an
economically meaningful cost, not a negligible one. We report our own figure as an illustrative order of
magnitude for a policy audience, not a causal welfare estimate.

The remainder of the paper is organized as follows. Section~\ref{sec:lit} reviews the two literatures this
paper connects. Section~\ref{sec:institutional} details the five states' entry-licensing institutions.
Section~\ref{sec:data} describes the data and the time-varying border construction. Section~\ref{sec:methods}
presents the empirical strategy. Section~\ref{sec:results} reports results. Section~\ref{sec:discussion}
discusses mechanism and limitations, and Section~\ref{sec:conclusion} concludes.

\section{Literature Review}
\label{sec:lit}

Two literatures bear directly on this paper's question, and, to our knowledge, neither has been connected to
the other.

The first documents that cannabis demand crosses state borders when legal status differs. \citet{HaoCowan2020}
show that recreational legalization draws substantial cross-border sales from residents of neighboring
nonlegal states. \citet{HansenMillerWeber2020} show the mechanism operating in the other direction: Washington
retailers near the Oregon border lost roughly 36 percent of their sales when Oregon's own retail market opened
in October 2015, evidence that the earlier premium those counties enjoyed was genuinely a captive-demand
effect. \citet{BaggioChongKwon2020} document a related cross-border effect between cannabis and a different
regulated good, showing that counties in states that legalized medical marijuana reduced alcohol purchases
relative to nearby counties across the state line, using a border-county-pair design in the tradition of
\citet{DubeLesterReich2010}. Outside cannabis, cross-border shopping in response to tax and regulatory
differences is a mature empirical literature; \citet{FribergEtAl2024} is a recent addition, using Norway's
COVID-19 border closures as a natural experiment to show that cross-border shopping measurably erodes domestic
commodity tax revenue.

The second literature studies how entry regulation shapes market structure and rents once a cannabis market
exists, independent of any cross-border demand story. \citet{HollenbeckUetake2021} document that Washington's
capped licensing regime supports meaningful markups. \citet{Thomas2019LicenseQuotas} builds a structural model
of Washington's own license quota and estimates that free entry would raise total surplus by roughly 18
percent relative to the actual capped and partly randomized allocation. This is, as far as we can establish,
the closest existing paper to the entry-regulation half of our own question -- but with no demand shock to
test it against. The general industrial-organization literature on entry supplies the theoretical scaffolding:
\citet{BresnahanReiss1990,BresnahanReiss1991} establish that competitive conduct changes most sharply with the
arrival of a market's second or third entrant; \citet{Berry1992}, \citet{Mazzeo2002}, and \citet{Seim2006}
extend this to settings with firm heterogeneity; \citet{ToivanenWaterson2005} shows a rival's existing
presence predicts a firm's own entry dynamically; \citet{BerryReiss2007} surveys the tradition. None of this
literature has been applied to cannabis border markets specifically, and none pairs an entry-regime moderator
with an exogenous demand shock of the kind cross-border legalization timing provides.

The gap is specific and modest, not a new mechanism. The cross-border demand literature has established that
demand moves across a cannabis-policy border; it has not asked where that demand goes once it arrives. The
entry-regulation literature has established that caps distort market structure and rents; it has not had an
exogenous demand shock to test that distortion against. This paper bridges those two literatures with a
design built for that purpose, and reports, rather than overstates, what the resulting test does and does not
show.

\section{Institutional Background}
\label{sec:institutional}

\begin{table}[htbp]
\centering
\caption{Sample states: legalization timing, border exposure, data, and entry regulation}
\label{tab:state_summary}
\footnotesize
\begin{tabular}{p{1.5cm}p{1.1cm}p{3.6cm}p{3.9cm}p{3.9cm}}
\toprule
State & Retail date & Nonlegal-neighbor exposure & Data source \& coverage & Entry-regulation regime \\
\midrule
Colorado & 2014-01-01 & KS, NE, OK, UT, WY (always); NM (nonlegal until 2022-04-01) & County-month sales, CO Dept. of Revenue (Socrata j7a3-jgd3), 2014--2025 & Capped (local-government opt-out bans; 48.4\% of counties ban all facilities, Pay\'an et al. 2021 -- see \S\ref{sec:institutional}) \\[0.4em]
Washington & 2014-07-08 & ID (always); OR (nonlegal until 2015-10-01) & County-fiscal-year sales, WA Liquor and Cannabis Board, FY2015--FY2025 & Capped (I-502 statewide license cap; 334$\to$556$\to$556+52 equity slots) \\[0.4em]
Oregon & 2015-10-01 & ID (always); CA (nonlegal until 2018-01-01), NV (nonlegal until 2017-07-01) & County-month sales, OLCC Socrata series, 2017--2025 (per-capita panel); 2025 cross-section (per-store, license registry) & Free 2015--2021, then application freeze (Jan 2022) and permanent per-capita cap (Mar 2024) \\[0.4em]
Montana & 2022-01-01 & ID, ND, SD, WY (always -- the only sample state bordering four nonlegal states) & County-month sales scraped from static DOR report pages, 2022--2025 (per-capita panel); 2025 cross-section (per-store, license list) & Capped (HB 701 local opt-in requirement for counties that voted against I-190 in Nov. 2020, binding from the 2022-01-01 launch; May-2025 moratorium separately restricts license issuance) \\[0.4em]
New Mexico & 2022-04-01 & OK, TX (always) & County-month sales AND licensee-level sales, NMRLD CROP portal, 2022--2026 (full panel, both outcomes) & Free (NMSA 1978 \S 26-2C-6(I) explicit statutory anti-cap provision; local bans preempted) \\[0.4em]
\bottomrule
\end{tabular}
\begin{minipage}{0.98\textwidth}
\vspace{0.3em}
\footnotesize
\textit{Notes:} ``Nonlegal-neighbor exposure'' lists every state sharing a land border, with the date any
neighbor's own retail-legalization date falls inside the sample state's panel window (i.e. requires a
time-varying, not static, border classification -- see \S\ref{sec:timevarying}). Entry-regulation
classifications are from statute-level research (Colorado Retail Marijuana Code / Amendment 64; Washington
RCW 69.50 / I-502; Oregon ORS 475C / HB 4016 / HB 4121; Montana HB 701 (MRTMA); New Mexico Cannabis
Regulation Act, NMSA 1978 Ch.\ 26 Art.\ 2C), not from labels.
\end{minipage}
\end{table}


Table~\ref{tab:state_summary} summarizes each sample state's legalization timing, border exposure, data
sources, and entry-regulation regime.

\textbf{Colorado} imposes no state-level numerical cap on retail marijuana licenses -- Amendment 64 (2012)
delegates licensing authority to local governments entirely. But \citet{PayanBrownSong2021}'s county-level
census of Colorado and Washington's local marijuana ordinances finds that, as of April 2019, 48.4 percent of
Colorado counties banned all marijuana facilities outright, and only 23.4 percent allowed every facility type
without restriction. We classify Colorado as capped for its entire panel on this basis: a state law that says
``no cap'' does not describe a free-entry market if local governments can and routinely do close most of the
state's geography to retail entirely.

\textbf{Washington} caps entry directly. Initiative 502 (2012) set an initial statewide ceiling of 334 retail
licenses; the Washington State Liquor and Cannabis Board raised it to 556 by emergency rule on January 6, 2016,
and a social-equity program began issuing up to 52 additional licenses in 2024. The cap level changes three
times over the sample window, but the state is never free-entry at any point in it.

\textbf{Oregon} is the one state whose regime genuinely switches mid-panel, and we treat it separately
(Section~\ref{sec:results}). Retail licensing was open from the market's October 2015 launch until a producer
moratorium (2018--2019) and a full application freeze effective January 2, 2022 (House Bill 4016, March 2022).
House Bill 4121 (March 2024) replaced the freeze with a permanent formula, one retail license per 7,500
residents aged 21 and over. We code Oregon free through 2021 and capped from 2022 onward.

\textbf{Montana}'s enabling statute (HB 701, retail effective January 1, 2022) imposes no numerical ceiling on
dispensary locations, but restricts new licenses to conversions of existing medical marijuana providers and
requires counties that voted against the state's Initiative 190 (November 2020) to affirmatively opt in before
a dispensary can locate there -- binding from the market's first day, not merely from a later moratorium
(Senate Bill 27, May 2025). Montana is coded capped for its entire panel on the opt-in requirement alone.

\textbf{New Mexico} is the cleanest free-entry state in the sample. The Cannabis Regulation Act (NMSA 1978
\S 26-2C-6(I)) explicitly bars the state's Cannabis Control Division from limiting license counts, and
licensing is administrative -- a 90-day issue-or-deny window, no lottery or merit-scoring step. Local
governments retain zoning authority but cannot ban cannabis retailers outright. New Mexico is coded free
throughout.

\section{Data}
\label{sec:data}

\subsection{Sample coverage}
\label{sec:samplecoverage}

The estimation sample pools county-year observations from five states -- Colorado, Washington, Oregon,
Montana, and New Mexico -- over 2015--2025, the exact window varying by state with each state's own
retail-launch and reporting history: Colorado and Washington from 2015, Oregon from 2017, and Montana and New
Mexico from 2022, the year each state's adult-use retail market opened. New Mexico's incomplete 2026 partial
year is dropped throughout. Two outcomes are estimated on two correspondingly different samples, because
Oregon's and Montana's revenue-per-store data exist only as a single 2025 cross-section rather than a full
panel (Section~\ref{sec:sources}): the revenue-per-capita sample has $n=1{,}283$ county-years across 178
counties, and the revenue-per-store and dispensary-density sample has $n=964$ county-years across 165
counties. Both samples exactly match the regression sample behind Table~\ref{tab:main_pooled}.

A county is classified as border-exposed if it shares a common land border with a nonlegal-state county
\emph{as of that observation's own year} (Section~\ref{sec:timevarying}) -- interior counties are every other
county of a sample state. Because this classification is time-varying rather than fixed, a given county can be
border in some sample years and interior in others; Table~\ref{tab:sample_composition} reports both the number
of counties that are ever classified as border or interior and the underlying county-year counts, together
with the same breakdown by entry regime. Fourteen counties change classification within the revenue-per-capita
sample -- five in Colorado (as New Mexico's 2022 legalization removes them from border status), six in
Washington (as Oregon's 2015 legalization does the same), and three in Oregon (bearing on its 2017
observations relative to California) -- a direct illustration of why the classification has to operate at
the county-year level, not the county level.

\begin{table}[htbp]
\centering
\caption{Sample composition, by estimation sample}
\label{tab:sample_composition}
\begin{tabular}{lcc}
\toprule
 & Revenue per capita & Revenue per store / density \\
\midrule
County-year observations & 1,283 & 964 \\
Unique counties & 178 & 165 \\
\addlinespace
Border counties (ever, in sample) & 54 & 48 \\
Interior counties (ever, in sample) & 138 & 128 \\
Counties that switch classification & 14 & 11 \\
\addlinespace
Border county-years & 280 & 219 \\
Interior county-years & 1,003 & 745 \\
\addlinespace
Capped-entry county-years & 1,024 & 840 \\
Free-entry county-years & 259 & 124 \\
\bottomrule
\end{tabular}
\begin{minipage}{0.92\textwidth}
\vspace{0.3em}
\footnotesize
\textit{Notes:} Both columns exclude New Mexico's incomplete 2026 partial year and match
Table~\ref{tab:main_pooled} exactly. ``Border (interior) counties, ever'' counts a county once if it is
ever classified as border (interior) in at least one sample year; because the border classification is
time-varying (Section~\ref{sec:timevarying}), a county can appear in both categories across different
years -- ``counties that switch classification'' reports how many do. Capped/free-entry county-years use
the state-year \texttt{Capped} indicator (Section~\ref{sec:institutional}).
\end{minipage}
\end{table}


\subsection{Revenue and dispensary-count sources}
\label{sec:sources}

Colorado's monthly county-level marijuana sales come from the Colorado Department of Revenue's Socrata series
\citep{COMarijuanaSalesByCounty} (2015--2025 in the estimation sample). Washington's come from the Liquor and
Cannabis Board's fiscal-year county sales and excise-tax reports \citep{WALCBSalesByCounty} (FY2015--FY2025).
Oregon's monthly county sales come from the Oregon Liquor and Cannabis Commission's own Socrata series
\citep{OLCCCannabisSalesData} (2017--2025); Oregon's revenue-per-store outcome exists only as a 2025
cross-section, since OLCC's license registry \citep{OLCCLicenseData} records a license's most recent renewal
date rather than its original opening date and so cannot support a historical per-store panel. Montana's
monthly county sales were scraped from the Department of Revenue's static report pages
\citep{MTRevenueCountySales}, since no bulk download exists; a county-year enters Montana's sample only if
all twelve months are disclosed; months withheld under the state's own dispensary-count confidentiality rule
are reported as missing, not zero, so a county-year with any suppressed month is dropped rather than
estimated from partial data. New Mexico's data come from the Cannabis Control Division's Cannabis Reporting
Online Portal \citep{NMCCDReporting}, which publishes monthly sales both by municipality and, uniquely among
the five states, by individual licensee, so New Mexico's revenue-per-store outcome (2022--2025 in this
sample) is a genuine historical panel, not a cross-section. Dispensary counts for Colorado come from a
geocoded license dataset, point-in-polygon assigned to counties; Washington's dispensary counts come directly
from WSLCB's own current retailer list \citep{WALCBCannabisRetailers2026}, which already assigns each license
to a county. County population and land area are 2023 vintage U.S. Census Bureau estimates
\citep{CensusPopEst2023}, held fixed across each county's panel years rather than updated annually -- we use a
single vintage throughout for consistency with the border-distance and land-area measures, which are
similarly time-invariant except where the border classification itself changes (Section~\ref{sec:timevarying}).

\subsection{Time-varying border classification}
\label{sec:timevarying}

Border exposure is defined by literal county-to-state-line geometric adjacency, not a distance threshold: a
county is border-exposed if its polygon touches the polygon of a nonlegal-state county. We compute this
directly from Census county boundary files, intersecting each sample-state county's polygon (buffered by
1,000 meters in a projected coordinate system, to bridge small gaps left by map simplification between
adjacent polygons -- some real bordering pairs, such as Pacific County, WA and Clatsop County, OR across the
Columbia River, are digitized roughly 400 meters apart in the underlying shapefile despite genuinely sharing a
border) against every other state's dissolved polygon, then applying each touching neighbor's own legal status
year by year rather than a single fixed snapshot -- incorrect for any county-year in which a neighboring
state's own status differed from that snapshot. Three cases in this sample are large enough to matter.
Colorado borders New Mexico (legal 2022-04-01); for Colorado county-years before that date, seven Colorado
counties (La Plata, Costilla, Montezuma, Baca, Las Animas, Conejos, and Archuleta) share a border with a New
Mexico county. Washington borders Oregon (legal 2015-10-01); eleven Washington counties, including populous
Clark County and Pacific County across the Columbia River from Portland and Astoria respectively, share a
border with Oregon, bearing on Washington's earliest fiscal-year observations. Oregon itself borders
California (legal 2018-01-01) and Nevada (legal 2017-07-01); five Oregon counties (Curry, Klamath, Josephine,
Lake, and Jackson) share a border with California, and three (Malheur, Lake, and Harney) with Nevada, bearing
on Oregon's first panel year, 2017.

Applying each neighbor's own retail-legalization date year by year reclassifies 39 of 1,283 per-capita
county-year observations and 36 of 964 revenue-per-store observations relative to a single fixed (2025)
legal-status snapshot, concentrated in the early years of Washington's and Colorado's panels and Oregon's
first year, and produces the fourteen within-sample classification switches reported in
Section~\ref{sec:samplecoverage}.

\subsection{Entry-regime classification}
\label{sec:entryregime}

We code a binary \texttt{Capped} indicator at the state-year level from each state's actual licensing statute
and the level of government that controls entry, not from whether a state imposes an explicit numerical
ceiling; the institutional basis for each state's coding is described in full in
Section~\ref{sec:institutional}. Colorado, Washington, and Montana are capped for their entire panel --
Colorado and Montana because local-government opt-out or opt-in authority functions as a binding restriction
even absent a state-level cap, Washington because of its direct license ceiling. New Mexico is free throughout,
the one state whose enabling statute affirmatively bars any numerical restriction on licenses. Oregon is the
only state whose classification changes within its own panel: free from its October 2015 launch through 2021,
then capped from 2022 onward, when the Oregon Liquor and Cannabis Commission's application freeze (effective
January 2, 2022, later replaced by a permanent per-capita licensing formula) stopped new entry. This switch is
coded at the state-year level -- every Oregon county-year from 2015--2021 is free and every Oregon county-year
from 2022 onward is capped -- rather than treating Oregon as permanently one regime or the other; it is this
within-state variation, not a cross-state comparison, that identifies the entry-regime interaction separately
within Oregon (Section~\ref{sec:methods}). Limitations of coding this indicator at the state-year rather than
county-year level are discussed in Section~\ref{sec:discussion}.

\subsection{Testing Washington's two-phase cap history}
\label{sec:localbarrier}

Washington's license cap has two distinct phases worth testing against the state-year \texttt{Capped}
indicator: an original 334-license ceiling (2014--2015), oversubscribed from launch by lottery, and a raised
556-license ceiling (from January 2016), never fully filled ($\sim$79 percent utilization). We tested directly
whether the second phase should be treated as effectively free entry (Section~\ref{sec:wa2016split}) and found
it should not: WSLCB closed the general retail application window permanently on March 31, 2016, and
Washington's post-2016 border counties do not behave like a genuinely free-entry benchmark (New Mexico) on
revenue per capita or dispensary density. Washington remains coded capped for its entire panel throughout this
paper.

\subsection{Panel structure}
\label{sec:panelstructure}

The unit of observation is the county-year for every state-outcome combination except one: Oregon's and
Montana's revenue-per-store and dispensary-density outcomes exist only as a single 2025 cross-section rather
than a panel (Section~\ref{sec:sources}), so for those two states specifically the unit of observation for
those two outcomes is the county. The geographic unit is the county as defined by the Census Bureau's cartographic county boundaries, and the
time frequency is annual (Washington's underlying source reports by fiscal year, July--June, aligned to
calendar year for pooling with the other four states). The panel is unbalanced: a county-year enters the
estimation sample only if that state's revenue source discloses non-suppressed sales for the full year
(Section~\ref{sec:sources}), and, for the revenue-per-store and dispensary-density outcomes specifically, only
if the county has a positive reported dispensary count that year. In practice this means a county enters the
panel in the year it opens its first dispensary and can drop back out in a later year if its revenue is
subsequently suppressed for confidentiality (Montana) or its per-store data does not extend to every sample
year (Oregon, Montana). Of Colorado's 42 sample counties, for example, only 25 report a complete eleven-year
run (2015--2025); the rest enter the panel partway through as new counties begin reporting a dispensary. The
resulting sample sizes -- 1,283 revenue-per-capita county-years and 964 revenue-per-store/density
county-years -- match Table~\ref{tab:main_pooled} exactly (Section~\ref{sec:samplecoverage}).

\subsection{Outcome variables}
\label{sec:outcomevars}

Three outcomes are estimated, each in logs. \textbf{Revenue per capita} is a county's total annual reported
marijuana sales divided by its 2023 Census population estimate. \textbf{Revenue per dispensary} is the same
annual sales total divided by the county's reported dispensary count for that year (Section~\ref{sec:sources});
because dispensary counts are only available for Colorado and Washington as a genuine annual panel, and only
as a single 2025 snapshot for Oregon and Montana, this outcome mixes panel and cross-sectional variation across
states, disclosed directly in the regression specification (Section~\ref{sec:methods}) rather than treated as a
uniform panel. \textbf{Dispensary density} is the county's dispensary count per 1,000 residents, using the same
population denominator as revenue per capita. All three are estimated in natural logs, with coefficients
converted to percentage effects as $(e^{\hat\beta}-1)\times 100$ throughout the paper; a small number of
county-years with zero dispensaries are necessarily absent from the revenue-per-store and density specs, since
the log of zero is undefined and a positive dispensary count is itself a condition for sample entry
(Section~\ref{sec:panelstructure}).

\subsection{Dispensary-level panel: New Mexico}
\label{sec:nmdispensary}

Every result above uses county-year aggregates. To test whether the county-level border premium survives at
the level of the individual seller, we separately construct a genuine dispensary-level panel for New Mexico --
the only sample state where individual-firm revenue is available historically. New Mexico's Cannabis Control
Division publishes 51 monthly licensee-level files (April 2022 onward), each listing every active licensee's
street address and reported sales. We identify a dispensary by its normalized street address rather than its
licensee name, since the same location can appear under different names across months (ownership changes,
rebrands), while different addresses under the same name are genuinely different stores. This yields 1{,}253
unique dispensary locations. Each dispensary's entry month is the first calendar month its address appears;
231 dispensaries (18.4 percent) are already present in the panel's first month (April 2022, the adult-use
launch date) and are flagged as left-censored (almost certainly pre-existing medical-program dispensaries,
whose true entry date is unobserved), while the remaining 1{,}022 are true new entrants whose opening month we
observe directly. County assignment reuses the same ZIP-to-county crosswalk already validated for the
county-level New Mexico panel; we do not build a separate address-level geocoding pipeline, since no
address-level coordinate source exists for New Mexico in our data (unlike Colorado, where a separately
assembled geocoded license file gives coordinates and opening dates but not store-level revenue). The
regression sample drops a small number of dispensary-years with zero reported revenue, leaving 1{,}249
dispensaries and 3{,}703 dispensary-year observations. Because New Mexico's border exposure and entry regime
are both time-invariant across the sample, this exercise cannot test the entry-regulation interaction; it
serves as a firm-level robustness check on the existence of the border premium itself.

\subsection{Summary statistics}
\label{sec:summarystats}

\begin{table}[htbp]
\centering
\caption{Summary statistics, estimation sample}
\label{tab:summary_stats}
\footnotesize
\resizebox{\textwidth}{!}{%
\begin{tabular}{lrrrrr}
\toprule
Variable & Mean & SD & Min & Max & $n$ \\
\midrule
\multicolumn{6}{l}{\textit{Panel A: Full sample}} \\
Revenue per capita (\$) & 336.14 & 559.13 & 0.00 & 7,986.19 & 1,283 \\
Revenue per dispensary (\$000s) & 1,587.8 & 1,237.2 & 0.0 & 7,461.5 & 964 \\
Dispensaries per 1,000 residents & 0.295 & 0.361 & 0.028 & 3.205 & 964 \\
County population & 160,333 & 310,126 & 624 & 2,271,380 & 964 \\
\addlinespace
Border county (shares a border with a nonlegal state) & 0.227 & 0.419 & 0 & 1 & 964 \\
Capped entry-regime state-year & 0.871 & 0.335 & 0 & 1 & 964 \\
\addlinespace
\multicolumn{6}{l}{\textit{Panel B: By entry regime (revenue-per-store sample)}} \\
\quad Revenue per dispensary (\$000s), Free-entry state-years & 631.1 & 316.9 & 0.0 & 1,463.0 & 124 \\
\quad Revenue per dispensary (\$000s), Capped-entry state-years & 1,729.0 & 1,259.8 & 26.1 & 7,461.5 & 840 \\
\quad Dispensaries per 1,000, Free-entry state-years & 0.492 & 0.525 & 0.051 & 3.205 & 124 \\
\quad Dispensaries per 1,000, Capped-entry state-years & 0.266 & 0.320 & 0.028 & 2.494 & 840 \\
\bottomrule
\end{tabular}
}
\begin{minipage}{0.95\textwidth}
\vspace{0.4em}
\footnotesize
\textit{Notes:} Panel A reports the pooled estimation sample: $n=1{,}283$ county-years for revenue per
capita (Colorado, Washington, Oregon, Montana, New Mexico, 2015--2025) and $n=964$ county-years for revenue
per dispensary and dispensary density (same states; Oregon and Montana contribute only their single 2025
cross-section to these two outcomes, Section~\ref{sec:data}). New Mexico's incomplete 2026 partial year is
excluded throughout, matching Table~\ref{tab:main_pooled}. Revenue per capita and revenue per dispensary are
annual dollars; dispensaries per 1,000 residents uses each county's contemporaneous population. The minimum
revenue-per-capita value (rounds to \$0.00) is a single sparsely populated county-year with negligible
reported sales, not a data-entry error. Border county is an indicator for sharing a common land border with a
nonlegal-state county as of that observation's own year (Section~\ref{sec:timevarying}); capped entry-regime
state-year is the state-year indicator defined in Section~\ref{sec:institutional}. Panel B splits the
revenue-per-store sample's two dollar/count outcomes by entry regime; the free-entry group ($n=124$) is New
Mexico's entire per-store panel (2022--2025) -- Oregon's own free period (2015--2021) predates its per-store
data, which exists only as a single 2025, already-capped cross-section (Section~\ref{sec:data}), so Oregon
contributes zero free-entry observations to this particular split.
\end{minipage}
\end{table}


Table~\ref{tab:summary_stats} reports summary statistics for the two estimation samples described above. Mean
revenue per capita is about 37 percent higher under capped-entry regimes than free-entry regimes in raw,
uncontrolled terms, and mean revenue per dispensary is about 2.7 times higher; both raw comparisons pool
capped and free county-years across every state and year rather than holding population, year, or border
status fixed, so they should be read only as descriptive context for Table~\ref{tab:main_pooled}'s controlled
estimates, not as evidence in their own right.

\subsection{Spatial distribution of dispensaries}
\label{sec:dispensarymap}

\begin{figure}[htbp]
\centering
\includegraphics[width=\textwidth]{figures/fig7_dispensary_map.png}
\caption{Dispensary density by county, 2025, the five sample states. Shading is dispensaries per 1,000
residents, the same variable used in Table~\ref{tab:main_pooled}'s dispensary-density specification; counties
shown in white have no dispensary recorded in the 2025 estimation sample (an unbalanced-panel county-year that
does not meet the sample-entry condition in Section~\ref{sec:panelstructure}, not necessarily zero dispensaries
in reality). Red outlines mark border counties (share a common land border with a nonlegal state as of 2025); black
outlines mark the five sample states. Small dark points show individual geocoded dispensary locations for
Colorado and Oregon only -- the two states with a usable point-level address source in our data
(Section~\ref{sec:sources}); Washington, Montana, and New Mexico dispensaries are represented only through the
county-level shading, since no reliable point-level location source exists for those three states in our
data (Washington's own geocoded file substantially undercounts its true 2025 retailer count, and Montana and
New Mexico were never geocoded to points at all).}
\label{fig:dispensary_map}
\end{figure}

Figure~\ref{fig:dispensary_map} maps this same dispensary-count variable spatially. The five states look
visibly different from one another: New Mexico's and Montana's border counties (red outline) include some of
each state's darkest-shaded, highest-density counties, while Colorado's border counties are, if anything,
lighter than its own interior, and Washington's and Oregon's dispensaries -- both point locations and county
shading alike -- concentrate around population centers (greater Seattle, the Portland metro area) well away
from any border county. The map is descriptive; we return to what distinguishes these states from one
another in Section~\ref{sec:results}.

\section{Empirical Strategy}
\label{sec:methods}

The main specification is
\[
y_{ist} = \alpha + \beta_1 \, \text{Border}_{ist} + \beta_2 \, \text{Capped}_{st} + \beta_3\,
\text{Border}_{ist} \times \text{Capped}_{st} + \gamma \ln(\text{population}_{ist}) + \delta_t + \eta_i + \varepsilon_{ist}
\]
estimated separately for log(revenue per dispensary store), log(revenue per capita), and log(dispensary count
per 1{,}000 residents). $\beta_3$ is the coefficient of interest, testing whether the border premium differs
by entry regime. State fixed effects $\eta_i$ are included for revenue per capita, where Oregon's within-state
regime switch identifies $\beta_3$ independent of the cross-state contrast, but omitted for revenue per store
and density: Oregon's per-store data is a single 2025 cross-section, so \texttt{Capped} has no within-state
variation for any state in that sample and becomes collinear with state fixed effects, which we verified
directly produces an unstable, uninterpretable coefficient. Standard errors are clustered by county throughout
(165--178 clusters), with HC1 standard errors reported alongside for comparison. Every specification is also
estimated separately by state, and Oregon's own free/capped switch is tested a third way, splitting its own
panel at the January 2022 freeze.

\section{Results}
\label{sec:results}

\begin{table}[htbp]
\centering
\caption{Main pooled specification: border $\times$ entry-regime interaction}
\label{tab:main_pooled}
\footnotesize
\resizebox{\textwidth}{!}{%
\begin{tabular}{lccccccc}
\toprule
& \multicolumn{2}{c}{Border} & \multicolumn{2}{c}{Capped} & \multicolumn{2}{c}{Border $\times$ Capped} & \\
Outcome & Coef. (\%) & $p$ (cl.) & Coef. (\%) & $p$ (cl.) & Coef. (\%) & $p$ (cl.) & $n$ \\
\midrule
Revenue per store & $+0.395$ ($+48.4\%$) & 0.181 & $+1.050$ ($+185.9\%$) & 0.000 & $+0.081$ ($+8.5\%$) & 0.799 & 964 \\
Revenue per capita & $+1.005$ ($+173.2\%$) & 0.001 & $+0.110$ ($+11.7\%$) & 0.242 & $-0.637$ ($-47.1\%$) & 0.048 & 1283 \\
Dispensary density & $+0.683$ ($+98.1\%$) & 0.000 & $-0.259$ ($-22.8\%$) & 0.082 & $-0.569$ ($-43.4\%$) & 0.009 & 964 \\
\bottomrule
\end{tabular}
}
\begin{minipage}{0.98\textwidth}
\vspace{0.3em}
\footnotesize
\textit{Notes:} Pooled across Colorado, Washington, Oregon, Montana, and New Mexico. Coefficients are OLS,
log-linear outcome; \% is $(e^{\hat\beta}-1)\times 100$. $p$ (cl.) is county-clustered. Revenue-per-store and
dispensary-density specs omit state fixed effects because \texttt{entry\_capped} has no within-state variation
in the per-store sample (Oregon's per-store data is a single 2025 cross-section), which would otherwise make
it collinear with the state dummies. Revenue-per-capita includes state FE since Oregon's per-capita panel
spans both its free (2015--2021) and capped (2022--2025) periods, giving genuine within-state identification.
Border is county-to-state-line geometric adjacency (shares a common land border with a nonlegal-state county),
time-varying by year (Section~\ref{sec:timevarying}), not a distance threshold. Entry-regime coding: Colorado,
Washington, and Montana are capped for their entire panel (local-government control over licensing -- county
opt-out bans in Colorado, opt-in requirements in Montana -- functions as a binding restriction even where no
state-level numerical cap exists); Oregon is free 2015--2021 and capped 2022--2025; New Mexico is free
throughout (statutory anti-cap provision, NMSA 1978 \S 26-2C-6(I)). The incomplete 2026 partial year is
excluded throughout.
\end{minipage}
\end{table}


Table~\ref{tab:main_pooled} reports the main pooled specification. The border premium is positive and
significant on two of three outcomes: revenue per capita is 173 percent higher ($p=0.001$), and dispensary
density is 98 percent higher ($p<0.001$); revenue per store is 48 percent higher but not statistically
distinguishable from zero under county clustering ($p=0.18$). The interaction with capped entry is a 47\%
relative reduction in the border premium's multiplier on revenue per capita (from roughly 173\% to roughly
44\% in level terms, $p=0.048$) and a 43\% relative reduction on dispensary density ($p=0.009$), but a precise
null on revenue per store (8.5 percent, $p=0.80$). Figure~\ref{fig:premium_by_regime} shows the pattern
directly.

\begin{figure}[htbp]
\centering
\includegraphics[width=0.85\textwidth]{figures/fig1_premium_by_regime.png}
\caption{Border premium by outcome and entry regime, with 95\% confidence intervals.}
\label{fig:premium_by_regime}
\end{figure}

\begin{table}[htbp]
\centering
\caption{State-specific border coefficients, by outcome}
\label{tab:state_specific}
\begin{tabular}{lccc}
\toprule
State (entry regime) & Revenue/store & Revenue/capita & Dispensary density \\
\midrule
CO (capped) & $+0.401$*** & $+0.270$*** & $-0.131$** \\
WA (capped) & $+0.469$*** & $+0.477$*** & $+0.008$ \\
OR (mixed) & $+1.168$*** & $+0.857$*** & $-0.522$ \\
MT (capped) & $-0.236$ & $+0.414$*** & $+0.259$** \\
NM (free) & $+0.365$* & $+0.989$*** & $+0.624$*** \\
\bottomrule
\end{tabular}
\begin{minipage}{0.85\textwidth}
\vspace{0.3em}
\footnotesize
\textit{Notes:} Each cell is the border-county coefficient from a separate within-state OLS regression
(log outcome on border indicator, log population, and year fixed effects where more than one year is
available), HC1 standard errors. $^{*}p<0.1$, $^{**}p<0.05$, $^{***}p<0.01$. ``--'' indicates no
border/interior variation within that state for that outcome.
\end{minipage}
\end{table}


Table~\ref{tab:state_specific} reports the same outcomes estimated separately within each state.
Colorado's own density coefficient is negative and significant ($-12$ percent) despite its large positive
revenue-per-store premium ($+49$ percent), matching the pooled interaction's prediction at the state level.
New Mexico's density coefficient is strongly positive ($+62$ percent). Montana's density coefficient is also
positive and significant ($+26$ percent) despite Montana being coded capped by its opt-in requirement -- the
opposite direction from Colorado's and the pooled interaction's own prediction, an anomaly we do not have a
confident explanation for and report directly rather than smooth over. Oregon's estimates are noisier,
reflecting small per-store and density subsamples (23 and 30 observations respectively).

\begin{figure}[htbp]
\centering
\includegraphics[width=0.85\textwidth]{figures/fig2_state_forest.png}
\caption{State-specific border coefficients by outcome, with 95\% confidence intervals (HC1).}
\label{fig:state_forest}
\end{figure}

\begin{table}[htbp]
\centering
\caption{Oregon within-state before/after: the 2022-01-01 entry freeze}
\label{tab:oregon_beforeafter}
\small
\begin{tabular}{lccc}
\toprule
Period & Border coefficient (\%) & $p$ (HC1) & $n$ \\
\midrule
Free (2017--2021) & $+1.004$ ($+172.9\%$) & 0.007 & 135 \\
Capped (2022--2025) & $+0.723$ ($+106.1\%$) & 0.101 & 120 \\
\midrule
\multicolumn{4}{l}{Formal interaction (border $\times$ post-freeze): $-0.381$ ($-31.7\%$), $p_{HC1}=0.503$,} \\
\multicolumn{4}{l}{$p_{cl}=0.611$, $n=255$} \\
\bottomrule
\end{tabular}
\begin{minipage}{0.85\textwidth}
\vspace{0.3em}
\footnotesize
\textit{Notes:} Revenue per capita, log-linear, county and log-population controls, year fixed effects.
Oregon's own before/after comparison points the same direction as the pooled cross-state interaction
(Table~\ref{tab:main_pooled}) -- the border effect shrinks and loses significance after the 2022 freeze --
but the formal within-Oregon interaction is not itself significant, consistent with the design being
underpowered on its own (only 6 Oregon border counties, adjacency definition). Revenue per store cannot be compared across the
freeze within Oregon: Oregon's per-store data exists only as a single 2025 (post-freeze) cross-section, so
there is no pre-2022 per-store observation -- reported as a limitation, not estimated.
\end{minipage}
\end{table}


Table~\ref{tab:oregon_beforeafter} and Figure~\ref{fig:oregon_beforeafter} split Oregon's own revenue-per-capita
panel at its January 2022 freeze. The border premium is larger in Oregon's free period (+173 percent,
$p=0.007$) than its capped period (+106 percent, $p=0.10$), the same direction as the pooled interaction,
though the formal within-Oregon interaction is not itself significant ($p=0.50$ HC1), which we attribute to
Oregon having only six border counties rather than to a contradiction of the pooled result.

\begin{figure}[htbp]
\centering
\includegraphics[width=0.7\textwidth]{figures/fig3_oregon_beforeafter.png}
\caption{Oregon's revenue-per-capita border premium, before/after the January 2022 freeze.}
\label{fig:oregon_beforeafter}
\end{figure}

\begin{figure}[htbp]
\centering
\includegraphics[width=0.72\textwidth]{figures/fig5_distance_gradient.png}
\caption{Revenue per store against log-distance to the nearest nonlegal-state county, residualized.}
\label{fig:distance_gradient}
\end{figure}

Replacing the binary border indicator with continuous log-distance recovers the same pattern (slope $-0.265$,
$p<0.001$; Figure~\ref{fig:distance_gradient}).

\begin{figure}[htbp]
\centering
\includegraphics[width=0.72\textwidth]{figures/fig6_captive_population.png}
\caption{Revenue per store against log nonlegal-side population within 70 miles, residualized.}
\label{fig:captive_population}
\end{figure}

The total population of nonlegal-state counties within 70 miles, the most robust mechanism evidence in our
earlier two-state analysis, weakens to a precise null pooling all five states (slope $+0.015$, HC1 $p=0.54$,
county-clustered $p=0.77$; Figure~\ref{fig:captive_population}) -- reported directly rather than retaining the
more favorable two-state framing.

\subsection{Robustness check: is Washington's cap genuinely binding?}
\label{sec:wa2016split}

Washington's license cap has two distinct phases the state-year \texttt{Capped} indicator does not capture:
an original 334-license ceiling (2014--2015), oversubscribed from launch by lottery, and a raised 556-license
ceiling (from January 2016). We tested directly whether this second phase should be treated as effectively
free entry, since the raised ceiling was never fully filled ($\sim$79 percent utilization). Two pieces of
evidence say no. First, WSLCB's general retail license application window closed permanently on March 31,
2016, and has not reopened since -- the unfilled balance is reserved for the Cannabis Social Equity Program
(only 16 of 52 dedicated slots issued as of the most recent count), not open competitive entry. Second, we
compared Washington's post-2016 border premium directly against New Mexico's, the paper's cleanest
genuinely free-entry benchmark.

\begin{table}[htbp]
\centering
\caption{Robustness check: Washington (2016 onward) vs.\ New Mexico, same specification as Table~\ref{tab:main_pooled}}
\label{tab:wa_vs_nm}
\footnotesize
\resizebox{\textwidth}{!}{%
\begin{tabular}{lccccccc}
\toprule
& \multicolumn{2}{c}{Border (NM level)} & \multicolumn{2}{c}{Washington (level shift)} & \multicolumn{2}{c}{Border $\times$ Washington} & \\
Outcome & Coef. (\%) & $p$ (cl.) & Coef. (\%) & $p$ (cl.) & Coef. (\%) & $p$ (cl.) & $n$ \\
\midrule
Revenue per store & $+0.357$ ($+42.9\%$) & 0.232 & $+0.999$ ($+171.5\%$) & 0.001 & $+0.071$ ($+7.3\%$) & 0.872 & 489 \\
Revenue per capita & $+0.967$ ($+162.9\%$) & 0.011 & $+0.212$ ($+23.6\%$) & 0.529 & $-0.483$ ($-38.3\%$) & 0.334 & 489 \\
Dispensary density & $+0.609$ ($+83.9\%$) & 0.000 & $-0.787$ ($-54.5\%$) & 0.000 & $-0.554$ ($-42.5\%$) & 0.007 & 489 \\
\bottomrule
\end{tabular}
}
\begin{minipage}{0.98\textwidth}
\vspace{0.3em}
\footnotesize
\textit{Notes:} Same specification and column structure as Table~\ref{tab:main_pooled} -- border $\times$
regime interaction, log(population) control, year fixed effects, county-clustered SEs -- with the regime
indicator replaced: \texttt{Washington} $=1$ for Washington restricted to 2016 onward, $=0$ for New Mexico
(the paper's cleanest genuinely free-entry benchmark, free for its entire panel; sample restricted to these
two states, no other states included). The border $\times$ Washington interaction tests whether Washington's
post-2016 border premium differs from New Mexico's; if Washington's post-2016 period genuinely behaved like
free entry, this interaction should be a precise null, matching the free-vs-capped null already established
for revenue per store elsewhere in this paper. On dispensary density it is not: the interaction is
significant and negative ($p=0.007$), the same direction as the main pooled interaction. On revenue per
capita the interaction points the same direction but is no longer significant at conventional levels under
county-clustered SEs ($p=0.334$) -- weaker than a naive HC1 read ($p=0.091$) would suggest, and we report the
weaker clustered result as the one to trust, consistent with this paper's practice throughout. Revenue per
store remains a genuine null ($p=0.872$). Taken together with the institutional record that WSLCB's
application window closed permanently in March 2016, the density result alone is enough to confirm
Washington's post-2016 period does not behave like New Mexico's free entry on the quantity margin, supporting
the paper's original classification (Washington capped for its entire panel) over the free-from-2016
alternative tested here.
\end{minipage}
\end{table}


Table~\ref{tab:wa_vs_nm} shows Washington's post-2016 period does not behave like New Mexico's on dispensary
density: the interaction is significant and negative ($-42.5$ percent, county-clustered $p=0.007$), the same
direction as the main pooled interaction. On revenue per capita the interaction points the same direction
($-38.3$ percent) but is not significant under county clustering ($p=0.334$, weaker than a naive HC1 read of
$p=0.091$ would suggest -- we report the more conservative clustered result). Revenue per store remains a
genuine null ($p=0.872$). The density result alone is enough to confirm the institutional record directly:
Washington continues to behave like a capped market on the quantity margin post-2016, so we retain the
paper's original classification (Washington capped throughout, Table~\ref{tab:main_pooled}) rather than the
free-from-2016 reclassification, which we tested and report here as a robustness check that only partially
held up.

\begin{figure}[htbp]
\centering
\includegraphics[width=\textwidth]{figures/fig_eventstudy_adjacency.png}
\caption{Border $\times$ year event study, all three outcomes, county-clustered 95\% confidence intervals.
Each point is the pooled border coefficient's deviation from the panel's first sample year (2015 for revenue
per capita and revenue per store), from a specification with border main effect, border $\times$ year
interactions, year fixed effects, and (for revenue per capita) state fixed effects. The dispensary-density
path grows over the panel, from roughly $+0.40$ log points in 2016 to $+0.63$ by 2023 -- the same direction as
the pooled interaction, not just an average of a flat relationship. Revenue per store shows the opposite, more
irregular pattern: its year-specific deviations are mostly negative and noisy (as low as $-0.52$ log points in
2022) rather than a steadily growing gap, which we report directly rather than force a clean trend narrative
onto the one outcome where the pooled interaction is itself a null.}
\label{fig:maincapped_eventstudy}
\end{figure}

\subsection{Firm-level robustness: a New Mexico dispensary panel}

\begin{table}[htbp]
\centering
\caption{New Mexico dispensary-level border premium, by year}
\label{tab:nm_dispensary_eventstudy}
\footnotesize
\begin{tabular}{lccc}
\toprule
Year & Border premium (\%) & Std.\ error (log pts) & $p$ (county-clustered) \\
\midrule
2022 & $+112.9\%$ & 0.322 & 0.019 \\
2023 & $+136.6\%$ & 0.285 & 0.002 \\
2024 & $+183.5\%$ & 0.301 & 0.001 \\
2025 & $+134.1\%$ & 0.304 & 0.005 \\
2026 & $+136.0\%$ & 0.275 & 0.002 \\
\midrule
Pooled (single coefficient, all years) & $+141.4\%$ & 0.272 & 0.001 \\
\bottomrule
\end{tabular}
\begin{minipage}{0.9\textwidth}
\vspace{0.3em}
\footnotesize
\textit{Notes:} $n=3{,}703$ dispensary-year observations, 1{,}249 unique dispensaries, 32 counties. Dependent
variable: log(annual dispensary revenue). Specification: log(revenue)$_{dt} = \alpha + \beta\,\text{Border}_d +
\sum_k \theta_k(\text{Border}_d \times \mathbf{1}\{\text{year}=k\}) + \gamma_1\ln(\text{county
population}) + \gamma_2\ln(1+\text{captive population}) + \delta_t + \varepsilon_{dt}$, county-clustered
standard errors. Border$_d$ is the home county's existing 70-mile nonlegal-border indicator (New Mexico's
border, to Texas and Arizona, is static across the sample -- both remain nonlegal throughout). Year-specific
premium = $\beta + \theta_k$ (2024 is the omitted reference year, so its row reports $\beta$ directly). The
pooled row is a separate, single-coefficient specification (\texttt{Border}, no year interactions) reported for
comparison. log(1+captive population) coefficient: $-0.067$ ($p<0.001$); log(county population) coefficient:
$-0.075$ ($p=0.16$).
\end{minipage}
\end{table}


Table~\ref{tab:nm_dispensary_eventstudy} and Figure~\ref{fig:nm_dispensary} re-estimate the border premium
directly on New Mexico's individual dispensaries (Section~\ref{sec:nmdispensary}) rather than on county
totals. This subsection's own dispensary-level county assignments still use the earlier 70-mile border
definition rather than the adjacency definition used everywhere else in this paper; New Mexico's border-county
set differs by exactly one county between the two definitions (Harding County is border under the 70-mile
threshold but not adjacency; San Juan County is the reverse), so we expect the difference in results to be
small, but we have not rebuilt the firm-level panel to confirm this and flag it as an open item rather than
silently assume it away. The firm-level premium is large, positive, and significant in every year (113 to 184
percent, all $p<0.02$); a pooled single-coefficient version recovers $+141.4$ percent ($p=0.001$), close to
New Mexico's county-level contribution to the main estimate. Border dispensaries are also disproportionately
new entrants (87.2 percent vs.\ 78.3 percent interior), though descriptively rather than causally, since New
Mexico's own regime never varies. This exercise cannot test the entry-regulation interaction, but confirms the
border premium is not a county-aggregation artifact in the one state where individual-firm revenue is
observable.

\begin{figure}[htbp]
\centering
\includegraphics[width=0.75\textwidth]{figures/fig13_nm_dispensary_eventstudy.png}
\caption{Border premium on individual-dispensary revenue, by year, New Mexico firm-level panel
($n=1{,}249$ dispensaries).}
\label{fig:nm_dispensary}
\end{figure}


\subsection{Economic magnitude: an illustrative revenue-gap accounting}
\label{sec:magnitude}

The interaction estimate implies a large difference in level terms, not just in the log coefficients
themselves: converting $\beta_1$ and $\beta_1+\beta_3$ separately, the free-entry border premium is
$+173.2$ percent while the capped-entry border premium is $+44.4$ percent, so a capped-state border county
would need $1.89\times$ its actual revenue to match what the free-entry premium implies for a county of its
size. Applying this ratio to actual 2025 border-county revenue in the three states where the cap is
unambiguously binding in our sample (Colorado, Washington, Montana; Table~\ref{tab:magnitude}) implies roughly
\$470 million in aggregate revenue, and roughly \$107 million in state tax revenue at each state's own
statutory cannabis tax rate (Colorado 15 percent, Washington 37 percent, Montana 20 percent), that these
states' border counties would need to gain to match the free-entry premium's implied level, on a single
year's data.

\begin{table}[htbp]
\centering
\caption{Illustrative revenue-gap accounting, capped-state border counties, 2025}
\label{tab:magnitude}
\small
\begin{tabular}{lccccc}
\toprule
State & Actual revenue & Counterfactual & Revenue gap & Tax rate & Tax gap \\
\midrule
Colorado & \$236.2M & \$446.8M & \$210.6M & 15\% & \$31.6M \\
Washington & \$154.9M & \$292.9M & \$138.1M & 37\% & \$51.1M \\
Montana & \$135.9M & \$257.1M & \$121.2M & 20\% & \$24.2M \\
\midrule
Total & \$527.0M & \$996.8M & \$469.8M & -- & \$106.9M \\
\bottomrule
\end{tabular}
\begin{minipage}{0.9\textwidth}
\vspace{0.3em}
\footnotesize
\textit{Notes:} Counterfactual = actual $\times$ 1.89, the ratio of $(1+\text{free premium})$ to
$(1+\text{capped premium})$ implied by the main pooled interaction (Table~\ref{tab:main_pooled}), with
Washington coded capped for its entire panel (Section~\ref{sec:wa2016split}) and the incomplete 2026 year
excluded. Montana's tax rate is its state adult-use cannabis excise tax; local-option taxes (up to an
additional 3\% in some Montana counties) and Colorado's separate wholesale excise tax are excluded.
\end{minipage}
\end{table}

We present this explicitly as an illustrative accounting exercise, not a causal welfare estimate, for three
reasons. First, it assumes external validity: that Colorado, Washington, and Montana's border counties would
earn revenue at New Mexico's and free-period Oregon's own rate if their entry regime changed, an extrapolation
the cross-state design does not directly identify, and Oregon's own within-state before/after test
(Section~\ref{sec:results}) points the same direction but is not independently significant. Second, the figure
is a revenue gap, not a welfare loss: it says nothing about whether the difference reflects genuine unmet
demand (a real economic loss, e.g.\ continued cross-border leakage) or a different distribution of similar
underlying demand, since higher revenue per store under a binding cap is also consistent with the market
power documented by \citet{HollenbeckUetake2021} for Washington's own licensing regime, and our data contain
no prices or quantities that would let us adjudicate between these. Third, it is a single-year (2025) snapshot,
not a cumulative or discounted figure. For context, our \$470 million revenue-gap estimate is considerably
larger than \citet{Thomas2019LicenseQuotas}'s structural estimate that free entry would have raised total
surplus by roughly 18 percent under Washington's own license quota -- an expected difference, since a revenue
gap includes normal costs and rents that a surplus calculation nets out, but the two estimates agree
qualitatively that Washington's entry cap imposes an economically meaningful cost, not a negligible one. We
report our own figure as an order of magnitude for a policy audience, not as a precise estimate.

\section{Discussion}
\label{sec:discussion}

The defensible claim here is narrower than the paper's opening hypothesis. Cross-border cannabis demand
produces a measurable premium in border-county markets on every outcome we observe. Entry regulation
moderates that premium on the quantity margin specifically -- store count and per-capita revenue -- not on
price at any single store. The original hypothesis predicted regulation would redirect the premium from
quantity into price; the data show regulation blocks the quantity response without redirecting anything, and
the price-margin premium exists regardless of regime.

Three limitations. First, the entry-regime variable is state-year, not county-year, though the institutional
detail motivating Colorado's and Montana's capped classification is itself local; a county-level restriction
measure is a natural extension we did not build here. Second, Oregon's own before/after test is directionally
consistent with the pooled result but not independently significant, and should be weighted as corroborating
rather than confirmatory. Third, the captive-population result reverses statistical significance once the
sample expands from two states to five states, a reminder that small-sample mechanism evidence does not
automatically generalize. Fourth, we directly tested whether Washington's own 2016 cap increase should change
its classification -- capped only 2014--2015, free from 2016 onward (Section~\ref{sec:wa2016split}) -- and
found it should not: Washington's post-2016 border counties do not behave like New Mexico's, the paper's
genuinely free-entry benchmark, on dispensary density, the outcome where the difference is significant under
county clustering (Table~\ref{tab:wa_vs_nm}); the revenue-per-capita difference points the same direction but
is not significant under clustering, weaker than an HC1-only read would suggest. This is consistent with
WSLCB's application window having closed permanently in March 2016. We report this as a robustness check that
only partially confirmed the main classification, not as a result we are silently discarding.
Set against these limitations, the New Mexico dispensary-level panel (Section~\ref{sec:nmdispensary}) is a
genuine strength: it confirms the border premium itself is not an artifact of county-level aggregation.

\section{Conclusion}
\label{sec:conclusion}

Uneven legalization creates real spatial demand spillovers in legal cannabis markets, defined here by literal
county-to-state-line adjacency rather than a distance threshold. Entry regulation -- coded from actual
statutes and local-government authority, not from a state's own description of itself -- determines whether
that premium flows into market expansion: capped-entry states see a substantially smaller premium on revenue
per capita (47 percent smaller in relative terms) and dispensary density (43 percent smaller), while
free-entry states' extra demand shows up as more stores. The premium on revenue at any single store does not
depend on regulatory regime, a real departure from the paper's starting hypothesis that we report rather than
obscure. The result connects two literatures -- cross-border cannabis demand and cannabis entry regulation --
that have so far developed independently.

\bibliographystyle{apalike}
\bibliography{references}

\end{document}
